\chapter{OOP-Bilanz: Was ist Objektorientierung?}
\renewcommand{\chaptertitle}{OOP-Bilanz: Was ist Objektorientierung?}

\lehead[]{\normalfont\sffamily\hspace*{-2.00cm}\textcolor{white}{\colorbox{lightblue}{\makebox[1.60cm][r]{\thechapter}}}\hspace{0.17cm}\textcolor{lightblue}{\chaptertitle}}
\rohead[]{\textcolor{lightblue}{\chaptertitle}\normalfont\sffamily\hspace*{0.17cm}\textcolor{white}{\colorbox{lightblue}{\makebox[1.60cm][l]{\thechapter}}}\hspace{-2.00cm}}
%\chead[]{}
\rehead[]{\textcolor{lightblue}{AvHG, Inf, My}}
\lohead[]{\textcolor{lightblue}{AvHG, Inf, My}}

In den letzten vier Kapiteln habt ihr Hunde, Raupen, Ampeln, Drachen und
Billardkugeln programmiert. Jetzt ist der richtige Moment für einen kurzen
Schritt zurück: \textbf{Was genau unterscheidet diese Art zu programmieren
eigentlich von dem, was wir am Anfang (z.\,B.\ beim Wald) gemacht haben -- und
wozu braucht ihr das später im Abitur?}


\section{Zwei Denkweisen im direkten Vergleich}

Um zu verstehen, was \emph{objektorientierte Programmierung} (OOP) ausmacht,
hilft ein Blick auf die Alternative, mit der wir in Kapitel~8 und~9 begonnen
haben -- die sogenannte \emph{prozedurale (imperative) Programmierung}:

\begin{center}
\bgroup
\def\arraystretch{1.25}
\begin{tabularx}{\textwidth}{|p{3.4cm}|X|X|}
\hline
\textbf{Merkmal} &
\textbf{Ohne OOP (prozedural)}\newline \small\textit{(z.\,B.\ Hund ohne Fachklasse)} &
\textbf{Mit OOP (objektorientiert)}\newline \small\textit{(z.\,B.\ Klasse \texttt{Hund} + Anwendung)} \\ \hline
\textbf{Wo liegen die Daten?} &
Verstreut als einzelne Variablen im Hauptprogramm (\lstinline|bobbyX|, \lstinline|bobbyY|, \lstinline|charlieX|, \ldots). &
Gebündelt als \textbf{Attribute} direkt im jeweiligen Objekt (\lstinline|bobby|, \lstinline|charlie|). \\ \hline
\textbf{Wo liegt der Code?} &
In langen Methoden, denen alle Daten einzeln als Parameter übergeben werden müssen. &
Als \textbf{Methoden} direkt bei den Daten in der passenden Fachklasse. \\ \hline
\textbf{Was passiert bei 50 Objekten?} &
Unübersichtliches Variablen-Chaos; hoher Copy-\&-Paste-Aufwand. &
Die Fachklasse wird \textbf{einmal} geschrieben; beliebig viele Objekte (\lstinline|new ...|) nutzen denselben Bauplan. \\ \hline
\textbf{Schutz vor Fehlern} &
Jede Stelle im Programm kann jede Variable versehentlich überschreiben. &
Durch \lstinline|private| und \emph{Getter/Setter} prüft das Objekt selbst, welche Änderungen erlaubt sind. \\ \hline
\end{tabularx}
\egroup
\end{center}

\medskip
\noindent \textbf{Merksatz:} In der Objektorientierung besteht ein Programm nicht
mehr aus einer langen Liste von Befehlen, sondern aus \textbf{eigenständigen
Software-Bausteinen (Objekten)}, die ihre eigenen Daten kennen und wissen, was
sie damit tun können.


\section{Die vier Grundprinzipien der OOP auf einen Blick}

Wenn in Klausuren oder im Abitur nach den Merkmalen der objektorientierten
Programmierung gefragt wird, sind vor allem diese vier Konzepte gemeint:

\begin{compactenum}[1.]
  \item \textbf{Abstraktion (Klassen als Baupläne):}\\
  Man konzentriert sich auf die für das Programm wesentlichen Eigenschaften und
  Fähigkeiten eines Gegenstands (z.\,B.\ Position, Farbe, \lstinline|zeichnen()|,
  \lstinline|bewegen()|) und fasst sie in einer \textbf{Klasse} zusammen.
  Unwichtige Details bleiben weg.

  \item \textbf{Datenkapselung (Geheimnisprinzip):}\\
  Daten (\emph{Attribute}) und zugehörige Operationen (\emph{Methoden}) bilden
  eine Einheit. Nach außen hin werden die Attribute mit \lstinline|private|
  versteckt; der Zugriff erfolgt kontrolliert über eine öffentliche
  Schnittstelle (\emph{Getter} und \emph{Setter}). Das Objekt funktioniert wie
  eine \textbf{Black Box}.

  \item \textbf{Vererbung (Code-Wiederverwendung):}\\
  Spezialisierte Klassen (\emph{Unterklassen}) erben mit \lstinline|extends|
  alle nicht-privaten Attribute und Methoden einer allgemeineren
  \emph{Oberklasse} (z.\,B.\ \texttt{GesichtDrachen} erbt von \texttt{Drachen})
  und erweitern oder überschreiben nur das, was neu hinzukommt.

  \item \textbf{Polymorphie (\glqq Vielgestaltigkeit\grqq ):}\\
  Eine Variable vom Typ der Oberklasse kann zur Laufzeit auf Objekte
  verschiedener Unterklassen verweisen (z.\,B.\ kann eine Variable vom Typ
  \texttt{Lampe} eine \texttt{QuadratLampe} oder eine \texttt{RahmenLampe}
  speichern). Beim Aufruf von \lstinline|zeichnen()| entscheidet Java
  automatisch, welche Variante ausgeführt wird.
\end{compactenum}


\section{Wozu braucht ihr das im weiteren Unterricht und im Abitur?}

Ab der Qualifikationsphase (Q1/Q2) programmiert ihr keine Hunde oder Drachen
mehr -- aber genau dieselben OOP-Mechanismen bilden das Fundament für alle
Abitur-Schwerpunktthemen:

\begin{compactitem}
  \item \textbf{Grafische Oberflächen (Swing):} Jedes Programmfenster erbt mit
  \lstinline|extends JFrame| von einer fertigen Java-Klasse; Buttons,
  Textfelder und Listen sind fertige Objekte, die ihr über ihre Methoden
  steuert.
  \item \textbf{Nebenläufigkeit (Threads):} Parallele Abläufe entstehen, indem
  ihr eigene Klassen von \lstinline|Thread| ableitet und deren Methode
  \lstinline|run()| überschreibt.
  \item \textbf{Datenbanken \& Client/Server:} Netzwerkverbindungen
  (\lstinline|Socket|), Datenströme (\lstinline|InputStreamReader|) und
  Datenbankabfragen (\lstinline|ResultSet|) sind allesamt Objekte, die ihr über
  ihre Methodenschnittstellen ansprecht -- genau wie ihr es in Kapitel~10.7
  gelernt habt.
\end{compactitem}
